\chapter{Configuration}

\label{ch:config}

This chapter answers what the daemon reads at start, which of two files wins when both set
the same key, what it refuses to read at all, and how to change a setting without a restart.
It also states where every file an installation writes actually lives, because the most
common configuration fault is not a wrong value --- it is a second daemon reading different
files from the one the operator is looking at.

\section{Two files, one effective config}

There is a global file and, optionally, a project file. The global file is the one an
operator edits for a host; the project file is the one a repository carries, and it may set
only the settings a repository has any business setting.

\begin{center}
\begin{minipage}{\linewidth}
\centering
\resizebox{\linewidth}{!}{%
\begin{tikzpicture}[node distance=8mm and 12mm]
  \node[b] (flag) {\texttt{-{}-config P}\\\scriptsize the flag};
  \node[b, below=of flag] (env) {\texttt{\$PII\_CONFIG}};
  \node[b, below=of env] (def) {\texttt{\$HOME/.pii/pii.toml}\\\scriptsize the default path};
  \node[b, right=of env, xshift=6mm] (proj) {\texttt{<workspace>/.pii/pii.toml}\\\scriptsize the project layer};
  \node[b, right=of proj, xshift=6mm] (eff) {effective config};
  \node[b, below=of eff] (defaults) {built-in defaults\\\scriptsize every key has one};
  \draw[f] (flag) -- (env);
  \draw[f] (env) -- (def);
  \draw[f] (env) -- (proj);
  \draw[sf] (defaults) -- (eff);
  \draw[f] (proj) -- (eff);
\end{tikzpicture}}
\captionof{figure}{Resolution order. The first of flag, \texttt{\$PII\_CONFIG} and the default path decides which global file is read; the project file is then layered over it, and every unset key falls through to a built-in default.}
\end{minipage}
\end{center}

The default path is \texttt{\$HOME/.pii/pii.toml}, and it is resolved \emph{before} the runtime
directory. That ordering is deliberate: \texttt{[core].runtime\_dir} lives in the file that
sets it, so a config path derived from the runtime directory would be circular.

A file that does not exist is not an error at the default path: the daemon starts with
built-in defaults and then refuses on the missing model. The same is true of a
\texttt{-{}-config} path that does not exist, which is worth knowing: a typo in the flag does
not say \enquote{no such file}, it says \enquote{no model configured} a moment later.

\section{The project layer, and what it may not say}

A workspace may carry \texttt{<workspace>/.pii/pii.toml}. It is read at start and on
\texttt{pii sys reload}, through one function, so the two can never disagree about which file
is read. A project file that exists but does not parse is fatal; it never silently falls back
to the global file.

The merge is not a general merge. A project value wins for three things only:

\begin{center}
\begin{tabular}{@{}>{\raggedright\arraybackslash}p{0.30\textwidth}>{\raggedright\arraybackslash}p{0.58\textwidth}@{}}
  \toprule
  \textbf{Project may set} & \textbf{Why it is safe} \\
  \midrule
  \texttt{[core].model} & the model a repository's work should use \\
  \texttt{[llama].host} & where that model is served, per repository \\
  \texttt{[[model\_tier]]} & the routing table, replaced whole rather than merged \\
  \bottomrule
\end{tabular}
\end{center}

Everything else that decides where the daemon writes or what it may reach is refused, and the
refusal names the setting:

\begin{lstlisting}[language=bash]
error: project config sets global-only settings: [net], grant_ceiling,
[core].runtime_dir (only the global ~/.pii/pii.toml may set them)
\end{lstlisting}

The list is short and closed: \texttt{[net]}, \texttt{grant\_ceiling},
\texttt{[core].runtime\_dir}, \texttt{[core].log\_file}, \texttt{[core].audit} and
\texttt{[core].observer}. The reason is asymmetric risk. A repository that could widen
\texttt{[net] grants} would let a checkout grant itself egress; a repository that could move
\texttt{runtime\_dir} would point a daemon at a store nobody is watching. A model id, by
contrast, is exactly the kind of thing that should differ per project.

\begin{center}
\begin{minipage}{\linewidth}
\centering
\resizebox{\linewidth}{!}{%
\begin{tikzpicture}[node distance=8mm and 16mm]
  \node[b] (global) {global \texttt{pii.toml}\\\scriptsize everything};
  \node[b, right=of global] (project) {project \texttt{pii.toml}\\\scriptsize model, llama host, tiers};
  \node[b, right=of project] (eff) {effective config};
  \node[b, below=of project, fill=red!6] (refused) {refused keys\\\scriptsize net, ceiling, dirs, log, audit, observer};
  \draw[f] (global) -- (eff);
  \draw[f] (project) -- (eff);
  \draw[bk] (refused) -- node[p, right=2pt, midway]{named} (project);
\end{tikzpicture}}
\captionof{figure}{The layer is a whitelist, not a patch. A project file is a set of opinions about models; anything that decides where bytes land belongs to the host.}
\end{minipage}
\end{center}

\section{The file is strict}

Every section is deserialised with unknown fields denied. A misspelled key is a load error
that names the key, not a setting that silently does nothing. This is a security property
rather than a style: \texttt{sandbx = "workspace"} must not start an unconfined daemon.

Two refusals are worth their own paragraph because they have been hit in practice:

\begin{itemize}
  \item \texttt{[core].ollama\_host} is refused by name. The ollama provider was retired; the
    key is not merely unknown, it produces a message that says to set \texttt{[llama] host}
    instead.
  \item Two \texttt{[[model\_tier]]} entries with the same name, or an entry with an empty
    name or model, fail validation before the daemon binds a socket.
\end{itemize}

\section{The sections}

\begin{longtable}{@{}>{\raggedright\arraybackslash}p{0.24\textwidth}>{\raggedright\arraybackslash}p{0.64\textwidth}@{}}
  \toprule
  \textbf{Key} & \textbf{What it decides} \\
  \midrule
  \endhead
  \texttt{grant\_ceiling} & capability names no graph or node may request beyond; empty means
    no ceiling. An unknown capability name is an error. This is the one setting a graph
    author cannot argue with. \\
  \texttt{[core]} & the host's own settings: model, directories, log, switches \\
  \texttt{[llama]} & where \texttt{llamacpp-server:} models dial; default
    \texttt{http://127.0.0.1:8080/v1} \\
  \texttt{[net]} & egress policy: \texttt{grants} (host:port patterns) and
    \texttt{allow\_private}. Default: no grants at all. \\
  \texttt{[[worker]]} & a supervised child: \texttt{name}, \texttt{program},
    \texttt{args}, \texttt{env}, \texttt{max\_in\_flight} \\
  \texttt{[[channel]]} & an inbound message source and the graph it starts per message \\
  \texttt{[[mqtt]]} & one broker connection and the graph each message starts \\
  \texttt{[hal]} & hardware backend (\texttt{sim} or \texttt{gpiod}) and the pins, I2C and
    SPI a graph may claim. Absent = the HAL is off and a \texttt{hal} node fails. \\
  \texttt{[wasm]} & plugin host: the one directory modules load from, plus per-call
    \texttt{timeout\_ms}, \texttt{max\_memory\_mib} and \texttt{fuel}. Absent = off. \\
  \texttt{[memory]} & \texttt{path}; empty means the store beside the main store \\
  \texttt{[memory\_graph]} & the bounds on the entity graph and the tier policy \\
  \texttt{[[cron]]} & a named timer entry: \texttt{name}, \texttt{graph},
    \texttt{schedule}, \texttt{input} \\
  \texttt{[[model\_tier]]} & a named model: \texttt{name}, \texttt{model},
    \texttt{tools}, \texttt{max\_tokens} \\
  \texttt{[evolve.routing]} & step to model, for the steps \texttt{localize},
    \texttt{patch}, \texttt{review}, \texttt{research}, \texttt{orchestrator}. A typo'd step
    is an error, not a silent no-route. \\
  \bottomrule
\end{longtable}

\subsection*{\texttt{[core]}}

\begin{longtable}{@{}>{\raggedright\arraybackslash}p{0.27\textwidth}>{\raggedright\arraybackslash}p{0.45\textwidth}>{\raggedright\arraybackslash}p{0.18\textwidth}@{}}
  \toprule
  \textbf{Key} & \textbf{Meaning} & \textbf{Default} \\
  \midrule
  \endhead
  \texttt{model} & the model spec runs use when a graph names none & none \\
  \texttt{runtime\_dir} & sockets, installed graphs, crashes & per-uid \\
  \texttt{durable\_dir} & channel, timer and other durable stores & state\texttt{/durable} \\
  \texttt{session\_runs\_cap} & runs remembered per session in memory & built-in cap \\
  \texttt{trajectory\_dir} & when set, each run appends a JSONL trajectory & off \\
  \texttt{trajectory\_max\_bytes} & aggregate bound on the trajectory directory & 512 MiB \\
  \texttt{trajectory\_max\_files} & aggregate file bound & 50{,}000 \\
  \texttt{learn\_interval\_secs} & the unattended learn pass interval & 86{,}400 \\
  \texttt{log\_file} & daemon events as JSON lines & off \\
  \texttt{log\_rotate\_bytes} & rotate before a line passes this size & 10{,}485{,}760 \\
  \texttt{log\_keep\_files} & rotated files kept beside the live one & 5 \\
  \texttt{observer} & \texttt{on}/\texttt{off}; measurement only & \texttt{on} \\
  \texttt{audit} & \texttt{on}/\texttt{off}; measurement only & \texttt{on} \\
  \bottomrule
\end{longtable}

Three keys read a zero in a way that is easy to guess wrong, so they are stated here rather
than left to the sample. \texttt{learn\_interval\_secs = 0} switches the unattended learn
pass off; the field is optional and its absence means every 86{,}400 seconds.
\texttt{trajectory\_max\_bytes} and \texttt{trajectory\_max\_files} treat a zero as
\enquote{unset} and fall back to the defaults, never as \enquote{unbounded}, because the
trajectory directory is always bounded. \texttt{log\_rotate\_bytes = 0} is different again:
it is passed through, and the sample documents it as \enquote{never rotate} --- which is a
legitimate choice for a log a supervisor already rotates, and a bad one for a host nobody
watches.

The two measurement switches deserve a warning that the file cannot give: switching
\texttt{audit} off stops denials reaching \texttt{audit.log}, and switching \texttt{observer}
off detaches tool-call, store and provider logging. Both announce themselves as warning
events at start, and both are documented as measurement-only. An installation that runs with
either off is not observably the same installation as one that runs with them on.

\subsection*{\texttt{[memory\_graph]}}

These bounds are also command-line flags on \texttt{pii memory}, so a one-off query can be
narrowed without touching the file.

\begin{center}
\begin{tabular}{@{}>{\raggedright\arraybackslash}p{0.26\textwidth}>{\raggedright\arraybackslash}p{0.16\textwidth}>{\raggedright\arraybackslash}p{0.46\textwidth}@{}}
  \toprule
  \textbf{Key} & \textbf{Default} & \textbf{Refusal when exceeded} \\
  \midrule
  \texttt{max\_entities} & 10{,}000 & \texttt{entity\_cap} \\
  \texttt{max\_relationships} & 50{,}000 & \texttt{relationship\_cap} \\
  \texttt{max\_provenance} & 32 & \texttt{provenance\_cap} \\
  \texttt{max\_attributes} & 32 & none; the extra attribute is refused \\
  \texttt{max\_depth} & 3 & depth clamp at the call \\
  \texttt{default\_depth} & 1 & --- \\
  \texttt{max\_results} & 200 & \texttt{result\_cap} \\
  \texttt{hot\_secs} & 86{,}400 & tier policy \\
  \texttt{warm\_secs} & 604{,}800 & tier policy \\
  \texttt{hot\_count} & 10 & tier policy \\
  \texttt{max\_bytes} & 16 MiB & \texttt{bytes\_cap} \\
  \bottomrule
\end{tabular}
\end{center}

\section{Where the files are}

Three directories, and they are not interchangeable. The state directory holds what must
survive; the runtime directory holds what a restart may rebuild; the workspace is what the
sandbox confines writes to.

\begin{center}
\begin{minipage}{\linewidth}
\centering
\resizebox{\linewidth}{!}{%
\begin{tikzpicture}[node distance=8mm and 12mm]
  \node[b] (state) {\texttt{\$HOME/.local/state/pii}\\\scriptsize state};
  \node[b, right=of state] (run) {\texttt{/run/user/<uid>/pii}\\\scriptsize runtime};
  \node[b, right=of run] (ws) {the workspace\\\scriptsize confinement root};
  \node[b, below=of state] (store) {\texttt{store.redb}};
  \node[b, below=of run] (sock) {\texttt{control.sock}\\\scriptsize \texttt{graphs/}};
  \node[b, below=of ws] (audit) {\texttt{audit.log}\\\scriptsize denials};
  \node[b, below=of store] (dur) {\texttt{durable/}\\\scriptsize channel, timer};
  \node[b, below=of sock] (sib) {\texttt{queue.redb}\\\scriptsize \texttt{memory.redb}};
  \draw[sf] (state) -- (store);
  \draw[sf] (run) -- (sock);
  \draw[sf] (ws) -- (audit);
  \draw[sf] (store) -- (dur);
  \draw[sf] (sock) -- (sib);
\end{tikzpicture}}
\captionof{figure}{The layout an operator has to hold in mind. The audit trail lives beside the store, not in the log directory, so it survives losing a log volume.}
\end{minipage}
\end{center}

\begin{center}
\begin{tabular}{@{}>{\raggedright\arraybackslash}p{0.30\textwidth}>{\raggedright\arraybackslash}p{0.58\textwidth}@{}}
  \toprule
  \textbf{Path} & \textbf{Contents and note} \\
  \midrule
  \texttt{\$PII\_STATE\_DIR} & the state directory; default
    \texttt{\$HOME/.local/state/pii} \\
  \texttt{\$PII\_RUNTIME\_DIR} & sockets and ephemeral files; default
    \texttt{/run/user/<uid>/pii} \\
  \texttt{store.redb} & the main store, in the runtime or state dir; \texttt{-{}-store}
    overrides \\
  \texttt{audit.log} & beside the store: denials and the sandbox profile in force, rotated
    like the log \\
  \texttt{queue.redb} & beside the store: the queue's items and events \\
  \texttt{memory.redb} & beside the store: documents and the entity graph \\
  \texttt{durable/channel.redb} & under the state dir: channel cursors and last started runs \\
  \texttt{durable/timer.redb} & under the state dir: the timer wheel \\
  \texttt{graphs/} & under the runtime dir: installed definitions, \texttt{<name>.toml} \\
  \texttt{./graphs} & the repository copy that \texttt{graph list} compares against \\
  \bottomrule
\end{tabular}
\end{center}

\texttt{pii graph list} prints \texttt{DRIFT} when the installed copy differs from the
repository copy, which is the fastest way to answer \enquote{is the daemon running what I
reviewed}.

\section{Secrets are not configuration}

No credential is a config value. The file names the \emph{variable} or the \emph{file} that
holds a secret, never the secret:

\begin{itemize}
  \item \texttt{[[channel]]} uses \texttt{token\_env} or \texttt{token\_file}, and
    \texttt{secret\_env} or \texttt{secret\_file} for webhook signatures.
  \item \texttt{[[mqtt]]} uses \texttt{username\_env} or \texttt{username\_file}, and
    \texttt{password\_env} or \texttt{password\_file}. Credentials in clear text to a
    non-loopback broker are refused unless \texttt{allow\_plaintext\_credentials} is true.
  \item Provider keys come from \texttt{OPENAI\_API\_KEY}, \texttt{ANTHROPIC\_API\_KEY} or
    \texttt{GEMINI\_API\_KEY} in the daemon's environment, and are scrubbed from provider
    errors and events.
\end{itemize}

The reason is not tidiness. A config file is copied into backups, printed in bug reports and
mounted into containers; an environment variable is not. The scrubber that removes
secret-named variables from log lines is built from the daemon's own environment, so a key
passed as a \texttt{-{}-config} value would be invisible to it.

\section{Environment variables}

\begin{longtable}{@{}>{\raggedright\arraybackslash}p{0.28\textwidth}>{\raggedright\arraybackslash}p{0.60\textwidth}@{}}
  \toprule
  \textbf{Variable} & \textbf{Effect} \\
  \midrule
  \endhead
  \texttt{PII\_CONFIG} & the config file, overriding the default path \\
  \texttt{PII\_STATE\_DIR} & the state directory \\
  \texttt{PII\_RUNTIME\_DIR} & the runtime directory; wins over
    \texttt{[core].runtime\_dir} \\
  \texttt{PII\_MODEL} & the model when \texttt{[core].model} is unset; the flag or the file
    wins \\
  \texttt{PII\_TRAJECTORY\_DIR} & trajectory output; overrides
    \texttt{[core].trajectory\_dir} \\
  \texttt{PII\_MQTT\_URL} & broker for \texttt{pii mqtt}, after \texttt{-{}-url} \\
  \texttt{PII\_ACTS} & the operator acts file the deploy-card and distance commands append to \\
  \texttt{PII\_LOOP\_EVENTS} & the loop event NDJSON file \\
  \texttt{PII\_AGENT} & the lane's agent: \texttt{claude} or \texttt{pi} \\
  \texttt{PII\_GRAPH\_EDIT\_OK} & permits graph edits from the spec lane \\
  \texttt{ALLOW\_DIRTY} & permits the lane to start on a dirty tree, for a resume \\
  \texttt{OPENAI\_BASE\_URL} & an OpenAI-compatible endpoint; any vendor with that shape \\
  \texttt{OPENAI\_API\_KEY} & its key; optional for a local server \\
  \texttt{ANTHROPIC\_BASE\_URL} & default \texttt{https://api.anthropic.com} \\
  \texttt{ANTHROPIC\_API\_KEY} & required for Anthropic models \\
  \texttt{GEMINI\_BASE\_URL} & the Google endpoint; the default is the public host \\
  \texttt{GEMINI\_API\_KEY} & sent as \texttt{x-goog-api-key}, never in the URL \\
  \bottomrule
\end{longtable}

\section{Changing a setting}

There are three cost levels, and choosing the wrong one is how a host ends up running a
config nobody edited.

\begin{center}
\begin{minipage}{\linewidth}
\centering
\resizebox{\linewidth}{!}{%
\begin{tikzpicture}[node distance=8mm and 12mm]
  \node[b] (edit) {edit the file};
  \node[b, right=of edit] (reload) {\texttt{pii sys reload}};
  \node[b, right=of reload] (status) {\texttt{pii health}};
  \node[b, below=of reload] (restart) {restart the daemon};
  \node[v, below=of status] (ok) {same config sha?};
  \draw[f] (edit) -- (reload);
  \draw[f] (reload) -- (status);
  \draw[f] (edit) |- (restart);
  \draw[f] (restart) -- (ok);
  \draw[f] (status) -- (ok);
\end{tikzpicture}}
\captionof{figure}{A change is not in force until a surface says so. \texttt{pii health} prints the config path and a sha256 of the file it read, which is the only proof that the daemon is running the file you edited.}
\end{minipage}
\end{center}

\texttt{pii health} prints, on its first three lines, the binary's version, target and
revision, this client's revision, and the config path with its digest:

\begin{lstlisting}[language=bash]
pii 0.2.0 (c05f3d1e4dd3, x86_64-unknown-linux-gnu)
client: c05f3d1e4dd3
config: /root/.pii/pii.toml sha256=9f2c...
\end{lstlisting}

If the digest is not the file you just saved, the daemon read a different file: check
\texttt{\$PII\_CONFIG} in the service unit before editing anything else.

\texttt{pii sys reload} is narrower than its name suggests, and knowing its exact scope saves
an afternoon. It re-reads the global file and the project layer, then applies exactly two
things: every configured worker is reloaded, and each running channel's graph name is
replaced. It returns the number of workers reloaded. Nothing else changes --- not the model,
not the memory bounds, not \texttt{[net] grants}, and not a channel's
\texttt{poll\_interval\_ms}, which is read once when the poll loop starts. For anything
outside those two, the daemon has to be restarted.

\begin{center}
\begin{tabular}{@{}>{\raggedright\arraybackslash}p{0.34\textwidth}>{\raggedright\arraybackslash}p{0.26\textwidth}>{\raggedright\arraybackslash}p{0.28\textwidth}@{}}
  \toprule
  \textbf{Setting} & \texttt{sys reload} & \textbf{Restart} \\
  \midrule
  \texttt{[[worker]]} entries & applied & applied \\
  \texttt{[[channel]].graph} & applied & applied \\
  \texttt{[[channel]].poll\_interval\_ms} & no & applied \\
  \texttt{[core].model}, \texttt{[llama].host} & no & applied \\
  \texttt{[net].grants}, \texttt{grant\_ceiling} & no & applied \\
  \texttt{[memory\_graph]} bounds & no & applied \\
  \texttt{[core].runtime\_dir}, \texttt{log\_file} & no & applied at next start \\
  \bottomrule
\end{tabular}
\end{center}

The boundary has one practical consequence worth stating: a policy change --- egress grants,
the capability ceiling, the model --- is never in force until a restart, so a host that has
been edited but not restarted is running a config nobody can see from the file.

\section{A production file, in full}

The repository ships \texttt{examples/pii.production.toml}, which lists every key with its
default. The parts an administrator changes first:

\begin{lstlisting}[language=bash]
[core]
model = "llamacpp-server:granite-4.2-3b"
runtime_dir = "/run/pii"
log_file = "/var/log/pii/pii.log"
log_rotate_bytes = 10485760
log_keep_files = 5

[llama]
host = "http://10.10.10.217:8082/v1"

[net]
grants = ["127.0.0.1:11434"]
allow_private = true

[memory]
path = ""

[[cron]]
name = "nightly"
graph = "night"
schedule = "0 3 * * *"
\end{lstlisting}

Two of those deserve comment. \texttt{[net] grants} is empty by default, so the file above
opens exactly one host:port pair and nothing else; a model served elsewhere needs its own
grant or every call is denied with \texttt{net.egress.denied}. And \texttt{schedule} is
display-only: the timer that actually fires is the graph's entry node, so a cron entry whose
graph does not start with a \texttt{timer} does not run on a clock at all.

\section{Refusals, and what they look like}

\begin{center}
\begin{tabular}{@{}>{\raggedright\arraybackslash}p{0.30\textwidth}>{\raggedright\arraybackslash}p{0.58\textwidth}@{}}
  \toprule
  \textbf{Symptom} & \textbf{Cause and check} \\
  \midrule
  \texttt{config: unknown field \ldots} & a misspelled key. The message names it; the file is
    strict on purpose. \\
  \texttt{[core].ollama\_host is retired} & the ollama provider is gone; use
    \texttt{[llama] host} with \texttt{llamacpp-server:<model>} \\
  \texttt{project config sets global-only settings} & the repository tried to set
    \texttt{[net]}, \texttt{grant\_ceiling}, a directory, the log or a switch \\
  \texttt{duplicate [[model\_tier]] name} & two tiers with one name; the routing table must be
    unambiguous \\
  \texttt{no model configured} (exit 5) & neither \texttt{[core].model} nor
    \texttt{\$PII\_MODEL}; also what a typo'd \texttt{-{}-config} path looks like \\
  exit 4 at start & the requested sandbox profile could not be enforced; the daemon refuses
    rather than run unconfined \\
  a setting that has no effect & the daemon read another file. Compare the
    \texttt{config: \ldots sha256=} line in \texttt{pii health} with the file on disk. \\
  \bottomrule
\end{tabular}
\end{center}

\begin{lstlisting}[language=bash]
pii daemon --config /etc/pii/pii.toml --workspace /projects/pi-improved \
           --sandbox workspace --log-file /var/log/pii/pii.log
pii health --json | jq '.config, .mounts'    # what it read, what it may touch
pii sys reload                               # re-read the file, no restart
\end{lstlisting}